\clearpage
\VerdierPhysicalPageAnchor{0172}%
\begingroup\small\setlength{\parskip}{0pt}\setlength{\emergencystretch}{1.5em}
\VerdierSemanticAnchor{ch4-4-4}
\medskip\noindent\textbf{4.4. Convergence.}

\VerdierSemanticAnchor{ch4-4-4-1}
\medskip\noindent\textbf{4.4.1.}\quad
Pour une étude détaillée de la convergence des suites spectrales, nous
renvoyons à \emph{loc. cit.} Nous n'utiliserons que le critère suivant :

\VerdierSemanticAnchor{ch4-4-4-2}
\medskip\noindent\textbf{Définition 4.4.2.}\quad
On dit qu'un objet spectral $(H^n(p,q),\delta)$ à valeurs dans une catégorie
abélienne $\mathcal A$ est \emph{stationnaire}, si pour tout entier $n$, il
existe deux entiers $u(n),v(n)\in\mathbb Z$ tels que :
\[
\begin{aligned}
H^n(p,q)&=0,&&u(n)\leq p\leq q\leq+\infty,\\
H^n(p,q)&=0,&&-\infty\leq p\leq q\leq v(n).
\end{aligned}
\]

\VerdierSemanticAnchor{ch4-4-4-3}
\medskip\noindent\textbf{Théorème 4.4.3.}\quad
\textit{Soit $(H^n(p,q),\delta)$ un objet spectral stationnaire.}

\textit{a) Pour tout $n$, la filtration
$F^pH^n(-\infty,+\infty)$, $p\in\widetilde{\mathbb Z}$, est stationnaire,
i.e. il existe deux entiers $n_1\leq n_2$ dans $\mathbb Z$ tels que :}
\[
F^{n_1}H^n(-\infty,+\infty)=0
\quad\text{et}\quad
F^{n_2}H^n(-\infty,+\infty)=H^n(-\infty,+\infty).
\]

\textit{b) Les différentielles partant de (resp. aboutissant à)
$E\!\left[\begin{smallmatrix}n\\r,p,q,s\end{smallmatrix}\right]$ sont nulles
si $r\leq v(n+1)$ (resp. $u(n-1)\leq s$), et de plus :}

\textit{c) Pour tout $n$ et tout couple $p\leq q$, les objets
$E\!\left[\begin{smallmatrix}n\\r,p,q,s\end{smallmatrix}\right]$ tels que}
\[
-\infty\leq r\leq v(n+1)
\quad\text{et}\quad
u(n-1)\leq s\leq+\infty
\]
\textit{sont canoniquement isomorphes.}

Ce théorème résulte immédiatement des définitions.

\VerdierSemanticAnchor{ch4-4-4-3-convergence}
Un objet spectral stationnaire fournit une suite spectrale fortement
convergente au sens de \emph{loc. cit.} Nous dirons qu'une suite spectrale
associée à un objet spectral est \emph{stationnaire}, ou encore,
abusivement, \emph{convergente}, lorsque cet objet spectral est stationnaire.
En utilisant l'une ou l'autre des notations de (4.3.3), nous écrirons :
\[
\mathrm I_r^{p,q}\Longrightarrow H^{p+q}(-\infty,+\infty),
\qquad
\mathrm{II}_r^{p,q}\Longrightarrow H^{p+q}(-\infty,+\infty)
\]
pour indiquer qu'une suite spectrale converge. L'ensemble des termes
$H^{p+q}(-\infty,+\infty)$ sera alors appelé l'\emph{aboutissement} de la
suite spectrale.
\endgroup
